\documentclass[10pt,a4paper]{article}

\usepackage[utf8]{inputenc}
\usepackage[T1]{fontenc}
\usepackage[french]{babel}
\usepackage[left=1.5cm,right=1.5cm,top=1.5cm,bottom=1.5cm]{geometry}
\usepackage{multirow,tabularx,booktabs,array}
\usepackage[table]{xcolor}
\usepackage{fouriernc}
\usepackage{amsmath}
\usepackage{fancyhdr}
\usepackage{colortbl}

\setlength{\extrarowheight}{1mm}
\renewcommand{\arraystretch}{1.35}
\setlength{\parindent}{0pt}

\pagestyle{fancy}
\renewcommand\headrulewidth{1pt}
\fancyhead[L]{Lycée Denis Diderot}
\fancyhead[R]{Première Bac Pro Microtechniques}
\fancyfoot[C]{}
\fancyfoot[R]{Année 2025-2026}
\fancyfoot[L]{Alexis RUIZ}

% Couleurs
\definecolor{exoheader}{RGB}{50, 100, 160}
\definecolor{exolight}{RGB}{210, 228, 248}
\definecolor{totalrow}{RGB}{255,230,200}
\definecolor{grandtotal}{RGB}{50,100,160}

% Macro en-tête de partie
\newcommand{\exotitle}[2]{%
  \noindent\colorbox{exoheader}{%
    \parbox{\dimexpr\linewidth-2\fboxsep\relax}{%
      \textcolor{white}{\textbf{\large #1 \hfill #2}}%
    }%
  }%
}

\begin{document}

\begin{center}
  {\LARGE\textbf{GRILLE DE CORRECTION}}\\[1mm]
  {\large Dissolution et Dilution -- Activité documentaire \hspace{1cm}
  Première Bac Pro Microtechniques \hspace{1cm} \textbf{Total : /20 points}}
\end{center}

\vspace{3mm}

% ===================== PARTIE A =====================
\exotitle{Partie A -- La dissolution}{/5 pts}

\vspace{1mm}

\begin{tabularx}{\linewidth}{|c|X|c|}
  \hline
  \rowcolor{exolight}
  \textbf{Q} & \textbf{Éléments de réponse attendus} & \textbf{Pts} \\
  \hline
  \multirow{2}{*}{1}
  & Soluté identifié : chlorure de sodium / NaCl
  & 0{,}5 \\
  \cline{2-3}
  & Solvant identifié : eau distillée
  & 0{,}5 \\
  \hline
  \multirow{2}{*}{2}
  & Formule correcte : $M = M(\text{Na}) + M(\text{Cl})$
  & 1 \\
  \cline{2-3}
  & \textbf{Résultat :} $M(\text{NaCl}) = 58{,}5$~g/mol (avec unité)
  & 1 \\
  \hline
  \multirow{2}{*}{3}
  & Conversion du volume : $V = 500~\text{mL} = 0{,}500~\text{L}$
  & 1 \\
  \cline{2-3}
  & \textbf{Résultat :} $C_m = 9$~g/L (avec unité)
  & 1 \\
  \hline
  \rowcolor{totalrow}
  \multicolumn{2}{|r|}{\textbf{Sous-total Partie A}} & \textbf{/5} \\
  \hline
\end{tabularx}

\vspace{4mm}

% ===================== PARTIE B =====================
\exotitle{Partie B -- La dilution}{/8 pts}

\vspace{1mm}

\begin{tabularx}{\linewidth}{|c|X|c|}
  \hline
  \rowcolor{exolight}
  \textbf{Q} & \textbf{Éléments de réponse attendus} & \textbf{Pts} \\
  \hline
  \multirow{2}{*}{4}
  & Définition correcte de la dilution (ajouter du solvant à une solution mère)
  & 0{,}5 \\
  \cline{2-3}
  & Ce qui est conservé : quantité de soluté / $C_1 V_1 = C_2 V_2$
  & 0{,}5 \\
  \hline
  \multirow{2}{*}{5}
  & Formule utilisée : $f = C_1 / C_2$
  & 1 \\
  \cline{2-3}
  & \textbf{Résultat :} $f = 8$ (sans unité)
  & 1 \\
  \hline
  \multirow{3}{*}{6}
  & Formule correcte et $V_1$ isolé : $V_1 = C_2 V_2 / C_1$
  & 1 \\
  \cline{2-3}
  & Calcul numérique correct
  & 1 \\
  \cline{2-3}
  & \textbf{Résultat :} $V_1 = 250$~mL (avec unité, conversion L→mL acceptée)
  & 1 \\
  \hline
  \multirow{2}{*}{7}
  & Pipette graduée pour prélever $V_1$ (justification : précision)
  & 0{,}5 \\
  \cline{2-3}
  & Fiole jaugée de 2,0~L pour préparer la solution fille (justification : volume exact)
  & 0{,}5 \\
  \hline
  \rowcolor{totalrow}
  \multicolumn{2}{|r|}{\textbf{Sous-total Partie B}} & \textbf{/8} \\
  \hline
\end{tabularx}

\vspace{4mm}

% ===================== PARTIE C =====================
\exotitle{Partie C -- Concentration molaire et bilan de synthèse}{/7 pts}

\vspace{1mm}

\begin{tabularx}{\linewidth}{|c|X|c|}
  \hline
  \rowcolor{exolight}
  \textbf{Q} & \textbf{Éléments de réponse attendus} & \textbf{Pts} \\
  \hline
  \multirow{2}{*}{8}
  & Formule correcte : $C = C_m / M$
  & 1 \\
  \cline{2-3}
  & \textbf{Résultat :} $C = 0{,}5$~mol/L (avec unité)
  & 1 \\
  \hline
  \multirow{2}{*}{9}
  & Formule correcte : $n = C \times V$
  & 1 \\
  \cline{2-3}
  & \textbf{Résultat :} $n = 2{,}5$~mol (avec unité)
  & 1 \\
  \hline
  10
  & Comparaison $C = 0{,}5 > 0{,}25$~mol/L et conclusion : bain encore utilisable
  & 1 \\
  \hline
  \multirow{2}{*}{11}
  & Différence dissolution/dilution correctement expliquée (point de départ solide vs solution)
  & 1 \\
  \cline{2-3}
  & Un exemple de chaque opération tiré des documents
  & 1 \\
  \hline
  \rowcolor{totalrow}
  \multicolumn{2}{|r|}{\textbf{Sous-total Partie C}} & \textbf{/7} \\
  \hline
\end{tabularx}

\vspace{6mm}

% ===================== TOTAL =====================
\noindent\colorbox{grandtotal}{%
  \parbox{\dimexpr\linewidth-2\fboxsep\relax}{%
    \textcolor{white}{\textbf{\Large \hfill TOTAL \hfill /20 points \hfill}}%
  }%
}

\vspace{6mm}

% ===================== TABLEAU ÉLÈVES =====================
\noindent\textbf{Récapitulatif de la classe :}

\vspace{2mm}

\begin{tabularx}{\linewidth}{|>{\centering\arraybackslash}p{4.5cm}|>{\centering\arraybackslash}X|>{\centering\arraybackslash}p{1.8cm}|>{\centering\arraybackslash}p{1.8cm}|>{\centering\arraybackslash}p{1.8cm}|>{\centering\arraybackslash}p{1.8cm}|}
  \hline
  \rowcolor{exolight}
  \textbf{Nom Prénom} & \textbf{Observations} & \textbf{Part. A /5} & \textbf{Part. B /8} & \textbf{Part. C /7} & \textbf{Total /20} \\
  \hline
  & & & & & \\[5mm] \hline
  & & & & & \\[5mm] \hline
  & & & & & \\[5mm] \hline
  & & & & & \\[5mm] \hline
  & & & & & \\[5mm] \hline
  & & & & & \\[5mm] \hline
  & & & & & \\[5mm] \hline
  & & & & & \\[5mm] \hline
  & & & & & \\[5mm] \hline
  & & & & & \\[5mm] \hline
  & & & & & \\[5mm] \hline
  & & & & & \\[5mm] \hline
  & & & & & \\[5mm] \hline
  & & & & & \\[5mm] \hline
  & & & & & \\[5mm] \hline
  & & & & & \\[5mm] \hline
  & & & & & \\[5mm] \hline
  & & & & & \\[5mm] \hline
  \rowcolor{totalrow}
  \textbf{Moyenne} & & & & & \\
  \hline
\end{tabularx}

\end{document}
